\SourceLeaf{011}pro omnibus valoribus ipsius \(x\) extra hos limites iacentibus semper sit \(=0\), manifesto etiam

\textit{valor integralis \(\int\varphi x.\mathrm{d}x\) ab \(x=-\infty\) usque ad \(x=+\infty\) extensi semper fit \(=1\).}

\Article{5}
Consideremus porro integrale \(\int x\varphi x.\mathrm{d}x\) inter eosdem limites, cuius valorem statuemus \(=k\). Si omnes errorum caussae simplices ita sunt comparatae, ut nulla adsit ratio, cur errorum aequalium sed signis oppositis affectorum alter facilius producatur quam alter, hoc etiam respectu erroris totalis valebit, sive erit \(\varphi(-x)=\varphi x\), et proin necessario \(k=0\). Hinc colligimus, quoties \(k\) non evanescat, sed e. g. sit quantitas positiva, necessario adesse debere unam alteramve errorum caussam, quae vel errores positivos tantum producere possit, vel certe positivos facilius quam negativos. Haecce quantitas \(k\), quae revera est medium omnium errorum possibilium, seu valor medius ipsius \(x\), commode dici potest erroris pars constans. Ceterum facile probari potest, partem constantem erroris totalis aequalem esse aggregato partium constantium, quas continent errores e singulis caussis simplicibus prodeuntes. Quodsi quantitas \(k\) nota supponitur, a quavis observatione resecatur, errorque observationis ita correctae designatur per \(x'\), ipsiusque probabilitas per \(\varphi' x'\), erit \(x'=x-k\), \(\varphi' x'=\varphi x\) ac proin \(\int x'\varphi' x'.\mathrm{d}x'=\int x\varphi x.\mathrm{d}x-\int k\varphi x.\mathrm{d}x=k-k=0\), i. e. errores observationum correctarum partem constantem non habebunt, quod et per se clarum est.

\Article{6}
Perinde ut integrale \(\int x\varphi x.\mathrm{d}x\), seu valor medius ipsius \(x\), erroris constantis vel absentiam vel praesentiam et magnitudinem docet, integrale
\[
\int xx\varphi x.\mathrm{d}x
\]
ab \(x=-\infty\) usque ad \(x=+\infty\) extensum (seu valor medius quadrati \(xx\)) aptissimum videtur ad incertitudinem observationum in genere definiendam et dimetiendam, ita ut e duobus observationum systematibus, quae quoad errorum facilitatem inter se differunt, eae praecisione praestare censeantur, in quibus integrale \(\int xx\varphi x.\mathrm{d}x\) valorem minorem obtinet. Quodsi quis hanc rationem pro arbitrio, nulla cogente necessitate, electam esse obiiciat, lubenter assentiemur.
